\chapter{AllSides Dataset Augmentation}
\label{sec:dataset}

This chapter presents our systematic extension of the AllSides political bias classification dataset \citep{baly2020we}. We collect additional articles from the same news sources present in the original dataset, holding the source set fixed so that we can evaluate whether the performance limitations observed in prior work were primarily attributable to data scarcity rather than inherent task difficulty.

AllSidesV2 adds 10,184 samples to the training split and 10,216 samples to the test split based on the sources seen in the original AllSides dataset. The validation set remained unchanged at 2,356 samples to ensure consistent evaluation metrics. This expansion represents approximately a $2.5\times$ increase in training data and a $8.5\times$ increase in test data compared to the original dataset. The distribution of political ideology labels remained roughly balanced across the extended samples. Appendix~\ref{appendix:extended_dataset} provides details about our data collection methodology and source-level statistics.  

\begin{table}[t!]
\centering
\small
\begin{tabular}{@{}lcccccccc@{}}
\toprule
\textbf{Models} & \multicolumn{4}{c}{\textbf{AllSidesV2 Media Split}} & \multicolumn{4}{c}{\textbf{AllSidesV2 Random Split}} \\
 & F1-Macro & Left & Center & Right & F1-Macro & Left & Center & Right \\
\midrule
\textbf{Transformer Baselines} & & & & & & & & \\
BERT & 40.95 & 65.20 & 20.11 & 37.53 & 81.90 & 84.63 & 79.90 & 81.16 \\
BART & 47.55 & 70.94 & 25.73 & 45.99 & 85.58 & 87.77 & 84.28 & 84.70 \\
RoBERTa & 47.83 & 70.66 & 25.73 & 54.86 & 87.83 & 89.91 & 85.11 & 88.46 \\
POLITICS & 60.60 & \textbf{77.94} & 38.55 & 65.30 & 88.23 & 90.60 & 84.95 & 89.12 \\
\midrule
\textbf{Fine-tuned LLM}$^\dagger$ & & & & & & & & \\
GPT-4o-mini & \cellcolor{blue!15}72.48 & \cellcolor{blue!15}81.52 & \cellcolor{blue!15}53.23 & \cellcolor{blue!15}82.69 & -- & -- & -- & -- \\
\midrule
\textbf{Pre-trained Transformers} & & & & & & & & \\
Best Full ($T_{id}(32){+}T_h$) & \textbf{62.19} & 75.97 & \textbf{39.83} & \textbf{70.76} & \textbf{89.29} & \textbf{90.91} & \textbf{86.48} & \textbf{90.47} \\
\bottomrule
\end{tabular}
\caption{Results on the AllSidesV2 dataset. Bold indicates best result per column among transformer and domain-adapted models. \colorbox{blue!15}{Highlighted} cells indicate where the fine-tuned LLM surpasses all transformer models.\\
$^\dagger$The fine-tuned LLM is not directly comparable to transformer models due to architectural differences and potential source leakage in the LLM pre-training corpus. This model was fine-tuned on 2000 samples for 3 epochs. More details can be found in Chapter~\ref{sec:llm-finetuning}.}
\label{tab:extended_main}
\end{table}

Re-evaluating transformer baselines on AllSidesV2 (Table~\ref{tab:extended_main}) reveals substantial but uneven gains: left F1 improves markedly for BERT and RoBERTa (12+ points), with more modest center gains, while right stagnates or declines across all models. Despite near-parity in left and right sample counts, the POLITICS model still shows a 12+ point F1 gap between the two classes (77.94 vs.\ 65.30), reinforcing that class balance alone does not resolve the challenge.

General-purpose transformers gain 2--12 points with additional data (with BART gaining the most), whereas POLITICS gains only 3 points, suggesting contrastive pre-training alone is approaching a ceiling. This motivates the specific design choice explored in Chapter~\ref{sec:continued-pretraining}: rather than scaling contrastive pre-training further, we add explicit metalinguistic auxiliary objectives (theme prediction and tone regression) on top of the contrastive backbone, so that the encoder is supervised on the structural cues that human annotators rely on rather than only on ideology-based positive/negative pairs.

\section{Evaluating Source Leakage}
\label{sec:source-leakage}

The two evaluation paradigms in Table~\ref{tab:extended_main} for our best-performing domain-adapted model ($T_{id}(32){+}T_h$, described in Chapter~\ref{sec:continued-pretraining}) differ by 27 percentage points in overall F1, where the media split provides the harder, more realistic evaluation. We decompose this gap by source, year, and topic. The gains from the source leakage is sharpest for Center-labeled articles: recall rises from 34.5\% in the media split to 80.3\% in the random split, with 78\% of Center errors in the media split predicted as Left. Per-source, per-year, and per-topic statistics for both paradigms are provided in Appendix~\ref{appendix:extended_dataset}.

The leakage benefit is highly non-uniform across outlets and falls into three tiers, as seen in Table~\ref{tab:source_leakage_tiers}. \emph{Tier~1} sources suffer the largest gap: their AllSides political label diverges systematically from their content register. Reuters (17.4\% $\to$ 33.0\%) and BBC News (18.5\% $\to$ 32.3\%) are rated Center by AllSides, yet the classifier predicts Left for their articles without source identity cues because their wire-service register resembles left-leaning editorial prose. At the other extreme, \emph{Tier~3} sources show minimal leakage benefit because their ideological register is content-identifiable: CNN (Web News) improves only marginally (33.0\% $\to$ 33.1\%), and Breitbart News similarly shows a minor uplift (31.4\% $\to$ 32.0\%). \emph{Tier~2} sources fall between these poles and highlight partisan outlets whose editorial style overlaps with an adjacent class, including Salon (26.4\% $\to$ 30.1\%) and ABC News (Online) (26.3\% $\to$ 32.8\%).

\begin{table}[h]                                                                                               
  \centering
  \small                                                                                                            
  \begin{tabular}{lr|lr|lr}                                                                                         
  \toprule                                                                                                          
  \multicolumn{2}{c|}{\textbf{Tier 1}} &                                                                          
  \multicolumn{2}{c|}{\textbf{Tier 2}} &                                                                            
  \multicolumn{2}{c}{\textbf{Tier 3}} \\                                                  
  \midrule                                                                                                          
  Source & $\Delta$ (pp) & Source & $\Delta$ (pp) & Source & $\Delta$ (pp) \\                                       
  \midrule                                                                                                          
  Reuters              & $+$15.6 & ABC News (Online) & $+$6.5 & CNN (Web News)        & $\approx$0 \\
  BBC News             & $+$13.8 & The Daily Caller  & $+$3.4 & American Spectator    & $+$0.1 \\
  ABC News             & $+$16.6 & Salon             & $+$3.7 & Guest Writer -- Left  & $+$0.2 \\
  Newsmax (News)       & $+$11.3 &                   &        & Guest Writer -- Right & $+$0.3 \\
  CBN                  & $+$9.1  &                   &        & Breitbart News        & $+$0.6 \\
  Newsmax              & $+$9.8  &                   &        & The Guardian          & $+$1.5 \\
  Associated Press     & $+$18.4 &                   &        & John Stossel          & $+$1.6 \\
  Reason               & $+$41.6 &                   &        & Guest Writer          & $+$3.4 \\
  \bottomrule
  \end{tabular}
    \caption{Per-source F1-macro gain from source leakage, grouped by tier. $\Delta$ = random-split F1-macro $-$ media-split F1-macro (percentage points). We show sources with $\geq 20$ articles in the test sets for both splits.}                                                         
  \label{tab:source_leakage_tiers}                                             
  \end{table} 

The contrast between Breitbart News and Newsmax illustrates that partisan label and content learnability are not equivalent. Both outlets are classified as right-hyperpartisan sources by the Ad Fontes Media Bias Chart\footnote{https://app.adfontesmedia.com/chart/interactive, retrieved 04/22/2026.}., yet their leakage profiles diverge sharply. Breitbart's distinctive rhetorical and lexical style produces reliable content-level signals regardless of whether the source appears during training (31.4\% vs.\ 32.0\%), placing it firmly in Tier~3. Newsmax, by contrast, combines partisan editorial commentary with heavy reliance on repurposed wire-service content and a tabloid-adjacent register, generating a systematic mismatch between its ``Right'' label and surface-level features; its F1-macro improves by approximately 10 percentage points with source leakage (18.5\% $\to$ 28.3\%). A parallel pattern appears on the left: Salon is rated Left by both AllSides and Ad Fontes, yet behaves as a Tier~2 source (26.4\% $\to$ 30.1\%) because its analytical prose style shares vocabulary with center-leaning outlets. These results suggest that editorial stance and content learnability are dissociable properties of a news source.

The temporal structure of the leakage effect provides further evidence that source routing, rather than content difficulty, drives the gap. In the media split, F1-macro rises from 48.3\% in 2012 to a peak of 62.1\% in 2015, then dips to 56.2\% in 2018 and settles in the 54--60\% range through 2024. The random split improves substantially over the same span, from 69.2\% in 2012 to over 90\% from 2017 onward, with only a modest dip in 2015--2016 (74.3\% and 71.4\%). The leakage gap consequently widens from approximately 12 percentage points in 2015--2016 to over 28 percentage points from 2017 onward, peaking at 34 percentage points in 2018. This widening coincides with the post-2016 shift toward meta-political and foreign policy coverage: reporting on trade policy, the Kavanaugh hearings, and the 2018 midterms was dominated by wire services such as Reuters and BBC, whose articles the classifier systematically misroutes to Left without source identity cues. The absence of any corresponding decline in the random split confirms this is a source-routing artifact rather than a shift in the inherent difficulty of political language.

Source-leakage dependence is most acute for content dominated by wire-service reporting. Topics covering disaster events (89.7\% miss rate in the media split), armed conflict (ukraine\_war: 66.7\%), and military affairs (us\_military: 56.5\%) improve dramatically when source identity is known, consistent with the finding that wire-service outlets are the primary Tier~1 sources. A qualitatively different set of topics remains difficult regardless of leakage, indicating genuine content-level ambiguity: race\_and\_racism (38.2\% miss rate in the random split) and polarization (34.8\%) involve articles that span ideological registers, while economic\_policy shows a persistent Center-to-Left confusion (100\% of Center errors go to Left in the random split), reflecting the substantial overlap between centrist and left-leaning economic vocabularies.

The error structure also shifts compositionally between splits. Extreme Left$\leftrightarrow$Right misclassifications rise from 25.3\% of all errors in the media split to 41.9\% in the random split (198 of 472 total errors). This is not a deterioration in performance but a shift in error composition: removing source leakage resolves the dominant Center$\leftrightarrow$Left confusion(45.8\% of media-split errors), thus exposing the harder cross-ideological cases that persist even when source identity is known. 